\scanpage{64}
\mspage{4.}
% archive foliation: 63
\diagram{\begin{tikzcd}
&X'\ \text{normalisé}\arrow[d]&\\
Y\arrow[r,phantom,"\subset" description]&X\text{--}\uline{\text{complet}}&
U=X-Y\\
&\text{réduit}&
\end{tikzcd}}
$Y_1\ldots Y_r$ les composantes de $Y$ de codim $1$ dans $X$

$Y'_{ij}$ les images inverses des $Y_i$ dans $X'$

\deleted{\unsure{Th}}
\deleted{\uline{Th} (1.)\quad$\unsure{\pi_1(U)}$
le quotient de $\pi_1(U)$ par le noyau de
$\pi_1(U)\longrightarrow\pi_1(X)$
est \unsure{un groupe de}}
Considérons les rev. étales de $X-Y$ qui
\deleted{\unsure{y sont définis}}
donnent des rev. \uuline{doux} des $\underline O_{\unsure{Y'_{ij}|Y_{ij}}}$,
ils sont classifiés par un groupe de type profini $\pi$.

\uline{Thm.} Le groupe $\pi$ est top. de génération finie.

\uline{Cor. 1} Soit $X$ \deleted{\unsure{Y}} complet,
$Y$ de codim $>1$, \inserted{$U=X-Y$,}
Alors $\pi_1(\unsure{U}\deleted{\unsure{X-Y}})$
est top. de génération finie.

\uline{Cor 2} Soit $X$ \deleted{de type} alg.
Alors $\pi'_1(X)$ est top. de type fini

\figure{\begin{tikzpicture}[x=1cm,y=1cm]
\draw (0,0)--(1.35,.08); \draw (.15,-.10)--(1.18,-.08);
\draw[->] (2.55,.35)--(2.55,.02);
\draw[->] (5,.27)--(5,-.27);
\draw (2.96,-2.95)--(3.22,-2.58);
\draw (3.02,-2.64)--(3.26,-2.96);
\end{tikzpicture}}
